\documentclass[10pt,a4paper]{amsart}
\usepackage[margin=19mm]{geometry}
\usepackage{amsmath,amssymb,mathtools}
\usepackage[scr=rsfs,cal=euler]{mathalfa}
\usepackage{xfrac}
\usepackage[unicode,hidelinks,bookmarks=false]{hyperref}
\newcommand{\ie}{i.e.\ }
\newcommand{\cf}{cf.\ }
\newcommand{\resp}{respectively\ }
\newcommand{\Subsection}{\subsection}
\providecommand{\nobreakdash}{\nobreak\hbox{-}}
\setlength{\emergencystretch}{2em}
\multlinegap0pt
\usepackage{amssymb}
\usepackage{xcolor}
\usepackage{enumerate}
\usepackage{dsfont}
\usepackage{float}
\newtheorem{proposition}{Proposition}
\newtheorem{lemma}{Lemma}
\newtheorem{definition}{Definition}
\newcommand{\D}{\mathbb{D}}
\newcommand{\N}{\mathbb{N}}
\newcommand{\QD}{\mathrm{QD}}
\newcommand{\R}{\mathbb{R}}
\newcommand{\SLE}{\mathrm{SLE}}
\newcommand{\cA}{\mathcal{A}}
\newcommand{\cL}{\mathcal{L}}
\newcommand{\clD}{\overline{\mathbb{D}}}
\newcommand{\ol}{\overline}
\newcommand{\qduration}{\mathcal{S}}
\title{Quantum Loewner evolution in quantum natural time: phases and Markov properties}
\author{Morris Ang, Deven Mithal}
\date{Source: arXiv:2605.03385v1 (CC BY 4.0)}
\begin{document}
\maketitle

\noindent\emph{Source and licence.} Quantum Loewner evolution in quantum natural time: phases and Markov properties. Version arXiv:2605.03385v1 (CC BY 4.0), distributed by arXiv under the Creative Commons Attribution 4.0 International licence, http://creativecommons.org/licenses/by/4.0/, as recorded in the arXiv licence metadata for this exact version and shown on the abstract page https://arxiv.org/abs/2605.03385v1. Full licence notice: \texttt{licenses/arxiv-2605.03385-CC-BY-4.0.txt}.\par\smallskip

\noindent\textit{Editorial scope.} Counted unit: the article's proposition prop-simple-markov (source0.tex lines 1141--1143). The article's complete original proof of it is delivered with it (source0.tex lines 1144--1175, one \texttt{proof} environment of 32 lines). No mathematics is written, repaired or re-derived: the statement and the proof are the article's own lines, in the article's own notation, and no retained line is substituted or re-flowed.  Retained uncounted as the source-local context the proof needs: lines 388--390; lines 395--397; lines 493--497; lines 514--517; lines 755--779; lines 795--797; lines 840--842; lines 1127--1129.  Nothing was omitted from any retained span, so the package carries no omission record.  No retained line contains a citation, so the wrapper carries no bibliography and no bibliography entry.  No retained line contains a figure environment, an image-inclusion command or drawing code, so no image is delivered and the package declares no asset dependency.  Editorial formatting only, in addition: an AMS \texttt{amsart} preamble that restates the article's own declarations and macros used by the retained lines, namely the environments \texttt{proposition}, \texttt{lemma}, \texttt{definition} on separate counters, together with the standard packages those lines need.  The retained environments print in this document as Proposition 1, Lemma 1, Definition 1 in order of appearance, whereas the article prints the same items with the numbering of its own full document.  The complete native source of the article as distributed in the arXiv e-print for this exact version is bundled unchanged as \texttt{sources/199786/source0.tex}.
\begin{proposition}\label{prop-conv-loc-spacetime}
Consider a convergent sequence of driving measures $(\nu^n)_{n \in \N}$ with limit $\nu$, and let $(g^n_t)_{t \geq 0}$ and $(g_t)_{t \geq 0}$ be the corresponding mapping out functions. Define $(g^n_\cdot)^{-1}: \D \times [0,\infty) \to \R$ via $(g^n_\cdot)^{-1}(z, t) := (g^n_t)^{-1}(z)$, and likewise define  $(g_\cdot)^{-1}$. As $n \to \infty$, we have $(g^n_\cdot)^{-1}\to (g_\cdot)^{-1}$ in the topology of uniform convergence on compact subsets of $\D \times [0,\infty)$. 
\end{proposition}

\begin{lemma}\label{lem-hausdorff-cara-conv}
Let $(K_n)_{n \in \N}$ be a sequence of compact hulls such that $K_n \to C$ in the Carath\'eodory topology and $K_n \to H$ in the Hausdorff topology. Then $\D \setminus C$ is the connected component of $\D \setminus H$ containing $0$.
\end{lemma}

\begin{definition}\label{def-quantum-time-simple}
    Let $\kappa < 4$ and $\gamma = \sqrt \kappa$. Let $h$ be a free boundary GFF on $\D$ and let $\eta: [0,\infty) \to \clD$ be an independent radial $\SLE_\kappa$ curve. For any $0 < t_1 < t_2$,
    the curve segment $\eta([t_1, t_2])$ corresponds to two boundary arcs of the simply-connected domain $\D \backslash \eta([0,\infty))$. 
    The \emph{quantum natural time} of $\eta$ elapsed between $t_1$ and $t_2$ is the $\gamma$-LQG boundary length in $(\D \backslash \eta([0,\infty)), h)/{\sim_\gamma}$ of either of these boundary arcs.
\end{definition}

\begin{definition}\label{def-quantum-time-sf}
    Let $\kappa >8$ and $\gamma = 4/\sqrt \kappa$. Let $h$ be a free boundary GFF on $\D$ and let $\eta: [0,\infty) \to \clD$ be an independent radial $\SLE_\kappa$ curve. For any $0 < t_1 < t_2$,
     the \emph{quantum natural time} of $\eta$ elapsed between $t_1$ and $t_2$ is $\cA^\gamma_h(\eta([t_1, t_2]))$.
\end{definition}

\begin{definition}[$\delta$-QLE$(\gamma^2, \eta)$ process] \label{def-delta-QLE-finite}
Consider either $\gamma \in (\sqrt3 -1,2)$ and $\eta = \frac3{\gamma^2} - \frac12$, or $\gamma \in (0, 4/3) \cup (2\sqrt3 - 2,2)$ and $\eta = \frac{3\gamma^2}{16} - \frac12$. In the former case set $\kappa = \gamma^2$, whereas in the latter case set $\kappa = 16/\gamma^2$.
Let $(\alpha,\beta)$ be defined as follows:
\begin{equation*} (\alpha, \beta) = 
        \begin{cases} 
          (Q-\frac{1}{\gamma},\gamma - \frac2\gamma) & \gamma \in (\sqrt3 - 1,2), \, \, \eta = \frac3{\gamma^2} - \frac12 \\
            (Q-\frac{\gamma}{4},\frac4\gamma - \frac\gamma2) & \gamma \in (2\sqrt3 - 2, 2), \, \, \eta = \frac{3\gamma^2}{16} - \frac12 \\
          (Q-\frac{\gamma}{4},\frac{3\gamma}{2}) & \gamma \in (0, 4/3), \, \,  \eta = \frac{3\gamma^2}{16} - \frac12 
       \end{cases}
    \end{equation*}
Sample $\phi^\delta \sim P_{\alpha, \beta, 1}$. We will define a shrinking process of origin-containing domains $D_s^\delta \subset \D$ beginning with $D_0^\delta = \D$. 
Iteratively for $n \in \mathbb Z_{\geq 0}$ we carry out the following:
    \begin{itemize}
        \item (\textbf{Resample the curve tip}): 
        Let $\phi^\delta_{n\delta} = g_{D^\delta_{n\delta}} \bullet_\gamma \phi^\delta$. 
        Sample a point from the probability measure $(\cL_{\phi_{n\delta}^\delta}^{\beta})^\#$ on $\partial \D$ and let $p^\delta_n \in \partial D_{n\delta}^\delta$ be its image under $(g_{D_{n\delta}^\delta})^{-1}$. Note that $p^\delta_n$ is understood as a prime end\footnote{If $\kappa \leq 4$, each point on the interior of an SLE$_\kappa$ segment corresponds to two prime ends in $\partial D^\delta_{n\delta}$, namely the boundary points on the left and right of the curve segment.} of $D_{n\delta}^\delta$. 
        
        \item (\textbf{Grow an $\SLE_\kappa$ curve}): Let $\eta^\delta_n$ be radial SLE$_\kappa$ in $D^\delta_{n\delta}$ from $p^\delta_n$ to $0$. 
        Parametrize $\eta_n^\delta$ by quantum natural time as in Definitions~
        \ref{def-quantum-time-simple}-\ref{def-quantum-time-sf}, and let $\qduration^\delta_n$ be its duration. 
        For $s \in (n \delta, n\delta + \min(\delta, \qduration^\delta_n)]$ let $D^\delta_s$ be the connected component of $D^\delta_{n\delta} \backslash \eta^\delta_n([0,s-n\delta])$ which contains 0. If  $\qduration^\delta_n \leq \delta$, set $\qduration^\delta = n\delta + \qduration^\delta_n$ and terminate the process at time  $\qduration^\delta$. 
    \end{itemize}
    If the process does not terminate, set $\qduration^\delta = \infty$. Let $K^\delta_s = \ol\D \backslash D^\delta_s$, so $(K^\delta_\cdot)_{[0,\qduration^\delta]}$ is a growing family of hulls. 
    We call $(\phi^\delta, (K^\delta_\cdot)_{[0,  \qduration^\delta]})$ a $\delta$-$\mathrm{QLE}(\gamma^2, \eta)$ process.    
\end{definition}

\begin{lemma}\label{lem-duration-indep} 
Defining $\phi_s^\delta := g_{K^\delta_s} \bullet_\gamma \phi^\delta$ and  $L_s^\delta := \cL^\gamma_{\phi_s^\delta}(\partial \D)$ for $s <\qduration^{\delta}$, we have $(\qduration^\delta,(L_\cdot^\delta)_{[0,\qduration^\delta)}) \stackrel d= (\qduration^\SLE, (L_\cdot^\SLE)_{[0,\qduration^\SLE)})$. In particular $\qduration^\delta$ is a.s.\ finite. Moreover, for any $s>0$, conditioned on $\{s < \qduration^{\delta}\}$  and on $(L^\delta_\cdot)_{[0,s]}$, the conditional law of $\phi_s^\delta$ is $P_{\alpha, \beta, L_s^\delta}$.
\end{lemma}

\begin{lemma}\label{lem-simple-unexplored-qle}
In the setting of Definition~\ref{def-delta-QLE-finite} and Lemma~\ref{lem-duration-indep}  with $\gamma \in (\sqrt3 - 1, 2)$ and $\eta = \frac{3}{\gamma^2} - \frac12$, the conditional law of $(\D \backslash  K^\delta_\qduration, \phi^\delta)/{\sim_\gamma}$ given $\qduration^\delta$ is $\QD(1 + 2 \qduration^\delta)^\#$.
\end{lemma}

\begin{lemma}\label{lem-simple-delta-massless}
    For each $k$ we have $\cA^\gamma_{\phi^{\delta_k}}(K^{\delta_k}_{\qduration^{\delta_k}}) = 0$ a.s.
\end{lemma}

\begin{proposition}\label{prop-simple-markov}
Passing to a further subsequential limit and recoupling using the Skorokhod representation theorem, we can guarantee the following. First, $K_s \subset K_\qduration$ for all $s < \qduration$ a.s. Second, $\cA^\gamma_\phi(K_\qduration) = 0$ a.s. Third, the conditional law of $(\D \backslash K_\qduration, \phi)/{\sim_\gamma}$ given $\qduration$ is $\QD(L_\qduration)^\#$.
\end{proposition}

\begin{proof}
\noindent
    For each $k$, sample a point $z^k \in \D$ from the probability measure $(\cA^\gamma_{\phi^{\delta_k}})^\#$ on $\D$; by Lemma~\ref{lem-simple-delta-massless} we have $z^k \in \D \backslash K^{\delta_k}_{\qduration^{\delta_k}}$ a.s. Let $f^k: \D \to \D \backslash K^{\delta_k}_{\qduration^{\delta_k}}$ be the conformal map such that $f^k(0) = z^k$ and $(f^k)'(0) > 0$. Independently sample $\theta$ uniformly from $[0,2\pi)$. 
    Passing to a subsequence and recoupling the laws via the Skorokhod representation theorem, we can ensure the limit $f = \lim_{k \to \infty} f^k$ a.s.\ exists in the topology of local uniform convergence. Since $\cA^\gamma_\phi$ has no atoms, we have $f'(0) \neq 0$. 
    \medskip 

    \noindent 
    \textbf{Step 1: The quantum surface parametrized by $f(\D)$ is a quantum disk.} 
   Let $R_\theta: \D \to \D$ be the rotation map $R_\theta(w) = e^{i\theta}w$. For each $k$ let $\psi^k = (R_{\theta} \circ (f^k)^{-1}) \bullet_\gamma \phi^{\delta_k}$ (so $(\D, \psi^k) \sim_\gamma (\D \backslash K^{\delta_k}_{\qduration^{\delta_k}}, \phi^{\delta_k})$), and let $\psi = (R_{\theta} \circ f^{-1}) \bullet_\gamma \phi$.

    By Lemma~\ref{lem-simple-unexplored-qle}, the conditional law of $(\D \backslash K^{\delta_k}_{\qduration^{\delta_k}}, \phi^{\delta_k})/{\sim_\gamma}$ given $\qduration^{\delta_k}$ is $\QD(1 + 2 \qduration^{\delta_k})^\#$, so since $\theta$ is independent of everything else, the conditional law of $\psi^k$ given $\qduration^{\delta_k}$ is $P(1 + 2 \qduration^{\delta_k})$ (defined immediately above this proof). 
    By Lemma~\ref{lem-duration-indep} the law of $\qduration^{\delta_k}$ does not depend on $k$, hence the joint law of $(\qduration^{\delta_k}, \psi^k)$ does not depend on $k$. Since $\qduration^{\delta_k} \to \qduration$ and $\psi^k \to \psi$ a.s., we conclude that $(\qduration, \psi) \stackrel d= (\qduration^{\delta_k}, \psi^k)$ for any $k$; in particular, the conditional law of $(\D, \psi)/{\sim_\gamma}$ given $\qduration$ is $\QD(1+2\qduration)^\#$. 

\medskip 
\noindent 
\textbf{Step 2: Showing $\cA^\gamma_\phi(\D \backslash f(\D)) = 0$ a.s.\ and $\cA^\gamma_\phi(K_s) = 0$ for rational $s < \qduration$ a.s.} For any fixed $k$, we have
    \[\cA^\gamma_{\phi^{\delta_k}}(\D) = \cA^\gamma_{\phi^{\delta_k}}(\D \backslash K^{\delta_k}_{\qduration^{\delta_k}}) = \cA^\gamma_{\psi^k}(\D) \stackrel d= \cA^\gamma_\psi(\D) = \cA^\gamma_{\phi}(f(\D)),\]
    where the first equality follows from Lemma~\ref{lem-simple-delta-massless}, the second from the definition of $\psi^k$, the distributional equality follows from $\psi^k \stackrel d= \psi$, and the last equality from the definition of $\psi$. Since $\phi^{\delta_k} \stackrel d= \phi$, we conclude that $\cA^\gamma_\phi(\D) \stackrel d= \cA^\gamma_\phi(f(\D))$, hence $\cA^\gamma_\phi(\D \backslash f(\D)) = 0$ a.s. By an analogous argument with Lemma~\ref{lem-simple-unexplored-qle} replaced by Lemma~\ref{lem-duration-indep}, for fixed $s$ the conditional laws of $\cA_\phi^\gamma(\D)$ and $\cA_\phi^\gamma (\D \backslash K_s)$ given $\{ s < \qduration\}$ agree, so the second claim follows. 

\medskip 
\noindent \textbf{Step 3: Showing $\D \backslash K_\qduration = f(\D)$.}
    %Now, note that a.s.\ the pointed domain $(\D \setminus f(\D), z)$ is the Carath\'eodory limit of $(K_{\qduration^k}^{\delta_k}, z)$. Indeed, 
    As $f^k \to f$ locally uniformly, %the point $z = f(0)$ satisfies $z \in f^k(\D)$ for all sufficiently large $k$, and so 
    the pointed domains $(f^k(\D), z)$ converge in the Carath\'eodory topology to $(f(\D), z)$. Since $f^k(\D) = \D \backslash K_{\qduration^{\delta_k}}^{\delta_k}$ and $\lim_{k \to \infty} K_{\qduration^{\delta_k}}^{\delta_k} = K_\qduration$ with respect to the Hausdorff topology, by Lemma~\ref{lem-hausdorff-cara-conv} the connected component of $\D \backslash K_\qduration$ containing $z$ is $f(\D)$. 
    If $\D \backslash K_\qduration$ contained a connected component other than $f(\D)$, then $\D \backslash f(\D)$ would have nonempty interior, contradicting Step 2 since $\cA^\gamma_\phi$ a.s.\ assigns positive mass to open sets. We conclude that $\D \backslash K_\qduration = f(\D)$. 

 
    Combining Steps 2 and 3 gives the second claim. By Step 3, we have $(\D \backslash K_\qduration, \phi)/{\sim_\gamma} = (f(\D), \phi)/{\sim_\gamma} = (\D, \psi)/{\sim_\gamma}$, and so by Step 1 the third claim holds.

  Finally, we turn to the first claim. By monotonicity it suffices to show that $K_s \subset K_\qduration$ for rational $s < \qduration$. Since $f^k \to f$ uniformly on compact sets and $f(\D) = \D \backslash K_\qduration$, for any point $w \in \D \backslash K_\qduration$ there exists a neighborhood $U \subset \D \backslash K_\qduration$ of $w$ such that $U \subset \D \setminus K_{\qduration^{\delta_k}}^{\delta_k}$ for all $k$ sufficiently large; we conclude that $U \subset \D \backslash \tilde K^{\delta_k}_{t^{\delta_k}(s)}$ for sufficiently large $k$.
    Since $\lim_{k\to\infty} t^{\delta_k}(s) = t(s)$, by Proposition~\ref{prop-conv-loc-spacetime} the Carath\'eodory limit of $\tilde K^{\delta_k}_{t^{\delta_k}(s)}$ is $\tilde K_{t(s)}$. Rewriting $\tilde K_{t(s)}$ as $K_s$, the Carath\'eodory kernel theorem yields that either $U \subset \D \backslash K_s$ or $U \subset K_s$; the latter is impossible since $\cA^\gamma_\phi(K_s) = 0$ by Step 2 but $\cA^\gamma_\phi(U) > 0$, hence $U \subset \D \backslash K_s$. Thus, any point $w$ in $\D \backslash K_\qduration$ is also in $\D \backslash K_s$, i.e., $K_s \subset K_\qduration$. This concludes the proof of the first claim. 
\end{proof}

\end{document}
